\chapter*{Реферат}
\thispagestyle{plain}

Пояснительная записка содержит 36 страницы, 16 рисунков. Количество использованных источников – 16.


Ключевые слова: сегментация, классификации, крюки и знамена, выборка, датасет, нейросети.


Целью работы является создание объемной и качественной выборки для нейросети, создание инструментов 
поддержки крюков и знамен в портале Корпуса.


В первом разделе приведен анализ методов сегментации и классификации крюков и знамен используемых в портале, 
подходов к формировании выборки для нейросетей, расположению крюков и знамен в древнеславянских текстах.


Во втором разделе приведена модель и алгоритмы для сегментации и классификации крюков.


В третьем разделе приведено описание архитектуры портала для реализации задачи сегментации и классификации.


В четвертом разделе представлена выборка и подход к ее формированию, для обучения нейросети классификации крюков.


%%% Local Variables:
%%% TeX-engine: xetex
%%% eval: (setq-local TeX-master (concat "../" (seq-find (-cut string-match ".*-3-pz\.tex$" <>) (directory-files ".."))))
%%% End:
